% Generated by analysis/word_lists_latex.py from preprocessing/05_remove_stopwords.py and
% config/boilerplate_phrases.txt; do not edit by hand. Requires array and longtable.

{\small
\renewcommand{\arraystretch}{1.3}
\setlength{\LTcapwidth}{\linewidth}
\begin{longtable}{|>{\raggedright\arraybackslash}p{0.24\linewidth}
                  |>{\raggedleft\arraybackslash}p{0.05\linewidth}
                  |>{\raggedright\arraybackslash}p{0.63\linewidth}|}

\caption{Stopwords (step 05): NLTK's Swedish list and the 86 words added for this corpus (200 stopwords in all). The added words were chosen by inspecting early models, with their lemma forms, since the text is lemmatized. Matched against the lemmatized, lowercased text.}
\label{tab:stopwords} \\
\hline
\textbf{Source} & \textbf{Words} & \textbf{Stopwords} \\
\hline
\endfirsthead
\hline
\textbf{Source} & \textbf{Words} & \textbf{Stopwords} \\
\hline
\endhead
\multicolumn{3}{r}{\emph{continued on next page}} \\
\endfoot
\endlastfoot

NLTK Swedish list & 114 & alla, allt, att, av, blev, bli, blir, blivit, de, dem, den, denna, deras, dess, dessa, det, detta, dig, din, dina, ditt, du, där, då, efter, ej, eller, en, er, era, ert, ett, från, för, ha, hade, han, hans, har, henne, hennes, hon, honom, hur, här, i, icke, ingen, inom, inte, jag, ju, kan, kunde, man, med, mellan, men, mig, min, mina, mitt, mot, mycket, ni, nu, när, någon, något, några, och, om, oss, på, samma, sedan, sig, sin, sina, sitta, själv, skulle, som, så, sådan, sådana, sådant, till, under, upp, ut, utan, vad, var, vara, varför, varit, varje, vars, vart, vem, vi, vid, vilka, vilkas, vilken, vilket, vår, våra, vårt, än, är, åt, över \\
\hline
Added for this corpus & 86 & all, alls, ar, bara, behöva, behöver, bland, bra, både, ca, del, dels, dock, ds, dvs, enligt, et, etc, exempel, fall, finna, finnas, fler, flera, flesta, frn, första, ge, genom, gjort, gjorts, gå, går, gör, göra, hel, hela, heller, hos, håll, hålla, håller, inför, iq, komma, kommer, kring, kunna, mer, mr, mången, ny, nya, of, olik, olika, osv, qz, resp, rätt, rör, röra, samt, se, ser, själva, ska, skall, sorn, stor, stora, sätta, sätts, såväl, ta, te, trots, ub, uö, vilja, vill, visa, visar, visat, viss, även \\
\hline
\end{longtable}
}

\clearpage

{\small
\renewcommand{\arraystretch}{1.3}
\setlength{\LTcapwidth}{\linewidth}
\begin{longtable}{|>{\raggedright\arraybackslash}p{0.24\linewidth}
                  |>{\raggedleft\arraybackslash}p{0.05\linewidth}
                  |>{\raggedright\arraybackslash}p{0.63\linewidth}|}

\caption{Parliamentary boilerplate (step 03): 450 words and phrases in 18 categories, matched against the original word forms; a word is removed only if a phrase covers the whole word.}
\label{tab:boilerplate} \\
\hline
\textbf{Category} & \textbf{Words} & \textbf{Words and phrases} \\
\hline
\endfirsthead
\hline
\textbf{Category} & \textbf{Words} & \textbf{Words and phrases} \\
\hline
\endhead
\multicolumn{3}{r}{\emph{continued on next page}} \\
\endfoot
\endlastfoot

OCR and Document Scanning Artifacts (these are errors, not content) & 45 & inskanna, inskannad, inskannade, inskannat, inskannar, inskanning, inskanningen, inscanna, inscannad, inscannade, inscannat, inscannar, inscanning, inscanningen, skanna, skannad, skannade, skannat, skannar, skanning, skanningen, scanna, scannad, scannade, scannat, scannar, scanning, scanningen, sidnummer, sidnumret, sidnumren, sidhuvud, sidhuvudet, sidhuvuden, sidfot, sidfoten, sidfötter, saml, sami, gotab, gotab stockholm, elanders, svenskt tryck, svenskt tryck stockholm, norstedts tryckeri \\
\hline
Government and Parliamentary Institutions & 27 & regering, regeringen, regeringens, regeringar, regeringarna, riksdag, riksdagen, riksdagens, riksdagar, riksdagarna, riksmötet, riksdagsåret, proposition, propositionen, propositionens, propositioner, propositionerna, motion, motionen, motionerna, motionär, motionären, motionärer, motionärerna, motion till riksdagen, utskottet, utskottets \\
\hline
Parliamentary Document and Reference Terms & 25 & hemställa, hemställan, hemställans, hemställer, hemställde, hemställt, hemställs, betänkande, betänkandet, observera, tabell, dokument, dokumentet, dokumentets, dokumenten, dokumentens, fel, yrkande, yrkandet, yrkanden, begär, begäran, förslaget, förslagen, tillkännagivande \\
\hline
Parliamentary Actions and Verbs & 48 & ställer, ställa, ställt, föreslå, föreslår, föreslog, föreslagit, anför, anförs, anföra, anföras, anfördes, anfört, anförda, anförts, anförande, anförandet, beslut, beslutet, beslutets, besluten, beslutens, beslutar, beslutade, hänvisning, hänvisa, hänvisar, hänvisade, hänvisat, hänvisas, hänvisades, anse, anser, ansåg, ansett, anses, ansågs, andel, andelen, andelar, andelarna, procent, procenten, procenttal, tillkännager, tillkännagivande, som sin mening, stockholm den \\
\hline
Question, Reference, and Opinion Words & 13 & fråga, frågan, vidare, mena, menar, skäl, skälet, skälen, borde, innebär, innebära, innebar, inneburit \\
\hline
Generic Filler and Discourse Words & 33 & vissa, stället, helt, sitt, tillägg, nuvarande, endast, utredning, känna, känner, kände, känt, kännas, kändes, kännedom, anledning, anledningen, anledningar, anledningarna, förekomma, förekommer, förekom, förekommit, förekomst, förekomsten, förekomster, eftersom, dessutom, möjligt, annat, annan, lydelse, övriga \\
\hline
Verb Forms and Grammar Words & 8 & sätt, satte, gäller, gällt, gällande, gällandes, ges, tar \\
\hline
Abbreviations and Single Characters (metadata/structural) & 5 & ant, åk, ab, ing, per \\
\hline
Time Period Names (temporal metadata) & 21 & dag, dagar, dagen, dagarna, vecka, veckor, veckan, veckorna, månad, månader, månaden, månaderna, år, året, åren, års, årtionde, årtionden, decennium, decennier, sekel \\
\hline
Weekdays (temporal metadata) & 7 & måndag, tisdag, onsdag, torsdag, fredag, lördag, söndag \\
\hline
Months (temporal metadata) & 12 & januari, februari, mars, april, maj, juni, juli, augusti, september, oktober, november, december \\
\hline
Seasons (temporal metadata) & 4 & vår, sommar, höst, vinter \\
\hline
Temporal and Frequency Expressions (metadata) & 43 & idag, i dag, igår, i går, imorgon, i morgon, tid, tiden, tider, period, perioden, perioder, nu, nyss, snart, sedan, sen, tidigare, senare, därefter, dessförinnan, fortfarande, redan, ännu, samtidigt, numera, hädanefter, framöver, hittills, länge, alltid, aldrig, ibland, ofta, sällan, dagligen, årligen, årsvis, månadsvis, veckovis, tillfälligt, slutligen, omedelbart \\
\hline
Spatial and Directional Words (structural metadata) & 53 & upp, uppåt, uppe, uppifrån, ned, ner, nedåt, nere, nedifrån, fram, framåt, framme, framifrån, bak, bakåt, bakom, bakifrån, in, inåt, inne, inifrån, ut, utåt, ute, utifrån, över, ovan, ovanpå, ovanför, under, underifrån, undertill, bredvid, intill, mellan, mitt, mitten, nära, bort, borta, bortåt, hit, dit, här, där, hem, hemma, hemåt, vart, varifrån, någonstans, överallt, ingenstans \\
\hline
Geographic Locations & 1 & sverige \\
\hline
Abbreviations and Parliamentary References (metadata/structural) & 21 & uhä, ex, prop, fl, mot, bl, fr, nr, kap, sid, sida, bil, dnr, sou, sfs, rskr, skr, bet, dir, prot, jfr \\
\hline
Currency and Numerical Terms (metadata) & 9 & kronor, miljoner, kr, milj, milj.kr, antal, antalet, antalen, antals \\
\hline
Political Party Names and Abbreviations (metadata from author signatures) & 75 & centerparti, centerpartiet, centerpartiets, centerpartister, centerpartisterna, socialdemokrat, socialdemokraten, socialdemokrater, socialdemokraterna, socialdemokratisk, socialdemokratiska, socialdemokratiskt, moderaterna, moderata samlingspartiet, sverigedemokrat, sverigedemokraten, sverigedemokrater, sverigedemokraterna, sverigedemokratisk, sverigedemokratiska, sverigedemokratiskt, folkparti, folkpartiet, folkpartiets, folkpartist, folkpartister, liberalerna, liberalernas, kristdemokrat, kristdemokraten, kristdemokrater, kristdemokraterna, kristdemokratisk, kristdemokratiska, kristdemokratiskt, miljöparti, miljöpartiet, miljöpartiets, miljöpartist, miljöpartister, vänsterparti, vänsterpartiet, vänsterpartiets, vänsterpartist, vänsterpartister, vänsterpartiet kommunisterna, ny demokrati, ny demokratis, s, m, sd, fp, mp, kd, vpk, nyd, c, fp, l, kd, kds, mp, v, vpk, nyd, ett, två, tre, fyra, fem, sex, sju, åtta, nio, tio \\
\hline
\end{longtable}
}
